\subsubsection{Trading Venues and Market Architecture}

\begin{principle} \hlt{Venue as Rule System}\\
A market is a rule system, not a place. Its rules govern access, order submission, visibility, execution, clearing, and settlement.
\end{principle}

\begin{definition} \hlt{Exchanges}\\
Exchanges centralise order interaction under membership, listing, trading, reporting, and clearing rules. They standardise access and make execution more reliable.
\end{definition}

\begin{definition} \hlt{OTC Markets}\\
OTC markets use bilateral or dealer-mediated negotiation rather than a central book. They fit bonds, swaps, FX, bespoke derivatives, and heterogeneous contracts.
\end{definition}

\begin{taxonomy} \hlt{Electronic Communication Networks (ECNs) and Alternative Trading Systems (ATSs)}\\
ECNs organise electronic matching outside traditional exchange structures. ATSs include dark pools, crossing systems, and specialised venues competing on anonymity, fees, block handling, or lower information leakage.
\end{taxonomy}

\begin{taxonomy} \hlt{Dark Pools, Internalizers, and Crossing Networks}\\
Dark pools hide pre-trade interest; internalizers match against broker-dealer inventory or internal flow; crossing networks match natural buyers and sellers, often at lit-market reference prices. These venues reduce signalling but can weaken public price discovery.
\end{taxonomy}

\begin{taxonomy} \hlt{Request for Quotation (RFQ) vs Central Limit Order Book (CLOB)}\\
RFQ markets request quotes from dealers, often bilaterally. A CLOB exposes standing interest multilaterally; RFQ protects relationships and discretion, while CLOBs improve comparability and public interaction.
\end{taxonomy}

\begin{caveat} \hlt{Stale Quotes and Vanishing Liquidity}\\
RFQ quotes can become stale before acceptance; displayed CLOB liquidity can vanish before an order arrives. In fragmented electronic markets, observed liquidity is an estimate, not a guarantee.
\end{caveat}

\begin{process} \hlt{Post-Trade Infrastructure}\\
Trading does not end at execution:
\begin{enumerate}[label=\roman*.]
\setlength{\itemsep}{0pt}
\item \hlt{clearing} matches buyer and seller records and validates trade terms;
\item \hlt{settlement} exchanges cash and securities;
\item \hlt{custody} holds assets on behalf of clients;
\item \hlt{depositories} immobilise or dematerialise securities so ownership can transfer by book entry.
\end{enumerate}
Clearinghouses may become buyer to every seller and seller to every buyer, reducing counterparty risk through margin and default management.
\end{process}

\begin{practitionernote} \hlt{Architecture Choice}\\
Venue choice depends on order size, urgency, information sensitivity, instrument standardisation, and counterparty risk. Small urgent orders favour visible and immediate venues; large or information-sensitive orders often favour dark, negotiated, or sliced execution.
\end{practitionernote}

\begin{mechanism} \hlt{Liquidity Bridges}\\
Fragmented venues require routing, arbitrage, and market-making links between books. These links reduce price dispersion, but their compensation is paid through spreads, fees, latency races, or information leakage.
\end{mechanism}

\begin{practitionernote} \hlt{Straight-Through Processing}\\
Straight-through processing sends trade details from execution to clearing, settlement, custody, and reporting with minimal manual intervention.
\end{practitionernote}
